\documentclass[12pt,a4paper]{article}
\usepackage[utf8]{inputenc}
\usepackage[spanish]{babel}
\usepackage{geometry}
\usepackage{hyperref}
\usepackage{booktabs}
\usepackage{listings}
\usepackage{xcolor}
\geometry{margin=2.5cm}

\definecolor{codegray}{rgb}{0.95,0.95,0.95}
\definecolor{codeblue}{rgb}{0.13,0.29,0.53}
\definecolor{codegreen}{rgb}{0.0,0.4,0.0}

\lstset{
  backgroundcolor=\color{codegray},
  basicstyle=\ttfamily\small,
  keywordstyle=\color{codeblue}\bfseries,
  commentstyle=\color{codegreen},
  stringstyle=\color{orange},
  breaklines=true,
  frame=single,
  numbers=left,
  numberstyle=\tiny\color{gray},
  showstringspaces=false,
  captionpos=b
}

\hypersetup{
  colorlinks=true,
  linkcolor=codeblue,
  urlcolor=codeblue,
  pdftitle={forge Framework -- Informe de Analisis Tecnico v2.0},
  pdfauthor={Analisis Adversarial Comparativo},
  pdfsubject={Agentic Development Framework}
}

\title{\textbf{forge Framework}\\
  \large Informe de An\'alisis T\'ecnico v2.0\\
  An\'alisis Adversarial Comparativo v1~$\rightarrow$~v2}
\author{An\'alisis Adversarial Independiente}
\date{3 de mayo de 2026}

\begin{document}

\maketitle
\tableofcontents
\newpage

% ----------------------------------------------------------------
\section{Resumen ejecutivo}

forge es un framework de agentic development que centraliza la configuraci\'on
de agentes de IA en un \'unico archivo \texttt{project.yaml} y genera, mediante
scripts determin\'isticos, los artefactos de configuraci\'on para m\'ultiples
runtimes (Claude Code, OpenCode, Kiro). Su propuesta arquitect\'onica ---taxonom\'ia
de tres tiers, compliance por dise\~no, Spec-Driven Development como workflow
obligatorio--- se mantiene intacta desde v1 y est\'a respaldada por evidencia de
c\'odigo.

La segunda ronda de an\'alisis (v2) eval\'ua el commit \texttt{d828157}, que
incorpora mejoras directamente motivadas por el an\'alisis v1. Las conclusiones
son mixtas. En las dimensiones que el an\'alisis v1 identific\'o como cr\'iticas,
la mejora es real y verificable: el bug de \texttt{install\_agent} est\'a
corregido, los adapters de OpenCode y Kiro tienen implementaciones funcionales,
el framework tiene teardown, el ecosistema de profiles pas\'o de cuatro a ocho, y
existe una suite de 112 tests. En las dimensiones donde los problemas son
estructurales ---coherencia entre documentaci\'on y c\'odigo, integraci\'on real
del flujo multi-runtime, ejecutabilidad de los fix messages del audit--- las
mejoras son parciales o nulas.

La recomendaci\'on consolidada es: forge v2 es adoptable para equipos de 3-8
personas que usen Claude Code como runtime principal, que entiendan que la
documentaci\'on central \texttt{agent-standard.md} no refleja el estado del
c\'odigo, y que no dependan de los fix messages del audit como comandos
ejecutables.

% ----------------------------------------------------------------
\section{Introducci\'on: contexto y metodolog\'ia}

\subsection{El framework}

forge se instala como git submodule en proyectos de desarrollo con agentes de IA.
Su funci\'on es proporcionar una base de agentes especializados, skills
componibles, y herramientas de auditor\'ia que hagan predecible y auditable el
comportamiento de los agentes en producci\'on.

La arquitectura central se basa en tres tiers:

\begin{itemize}
  \item \textbf{Tier 1 (Universal):} agentes que aplican a cualquier proyecto
    independientemente del stack.
  \item \textbf{Tier 2 (Profile):} agentes especializados en un stack t\'ecnico
    espec\'ifico (FastAPI, Rails, NestJS, etc.).
  \item \textbf{Tier 3 (Dominio):} agentes creados por el equipo para necesidades
    espec\'ificas del proyecto.
\end{itemize}

El \texttt{project.yaml} act\'ua como fuente de verdad: stack, roster de agentes,
frameworks de compliance, fases del sprint. Scripts determin\'isticos generan los
artefactos de configuraci\'on para cada runtime desde esa \'unica fuente.

\subsection{Metodolog\'ia}

Este informe consolida dos an\'alisis adversariales independientes realizados
sobre el mismo commit:

\begin{itemize}
  \item \textbf{An\'alisis favorable (\texttt{advocate-report.md} v2):} eval\'ua
    el dise\~no, las fortalezas, y las mejoras concretas incorporadas desde v1.
  \item \textbf{An\'alisis cr\'itico (\texttt{critic-report.md} v2):} examina la
    implementaci\'on en busca de problemas que afecten la adopci\'on en
    producci\'on.
\end{itemize}

Ambos an\'alisis se basan en lectura directa del c\'odigo fuente: scripts Python,
agentes Markdown, tests, hooks, adapters, y documentaci\'on. No se ejecut\'o el
framework en un proyecto real.

% ----------------------------------------------------------------
\section{Problemas del v1 resueltos en v2}

\subsection{Bug de \texttt{install\_agent} --- resuelto}

El an\'alisis v1 identific\'o que la funci\'on \texttt{install\_agent()} en
\texttt{forge-init.py} retornaba siempre \texttt{"UPDATE"} para instalaciones
nuevas, porque evaluaba \texttt{dst.exists()} despu\'es de copiar el archivo. La
correcci\'on consiste en capturar el estado del archivo antes de la operaci\'on de
copia:

\begin{lstlisting}[language=Python, caption={Correcci\'on en forge-init.py (l\'ineas 95--105)}]
# evaluar ANTES de copiar
already_existed = dst.exists()
shutil.copy2(src, dst)
return "UPDATE" if already_existed else "OK"
\end{lstlisting}

El test \texttt{test\_ok\_no\_es\_update\_en\_instalacion\_nueva} documenta
expl\'icitamente el bug original en su docstring. Correcci\'on completa.

\subsection{Flag \texttt{-{}-only} inexistente --- resuelto}

El audit generaba fix messages con \texttt{-{}-force -{}-only=<agente>}, pero el
flag \texttt{-{}-only} no exist\'ia. Ahora est\'a implementado en ambas sintaxis
(\texttt{-{}-only=nombre} y \texttt{-{}-only nombre}) con cobertura de tests en
dos archivos distintos.

\subsection{Adapters de OpenCode y Kiro vac\'ios --- parcialmente resuelto}

Los directorios \texttt{adapters/opencode/} y \texttt{adapters/kiro/} estaban
vac\'ios en v1. Ahora contienen implementaciones funcionales con 24 tests de
integraci\'on. La calificaci\'on es parcialmente resuelto porque el flag
\texttt{-{}-tool opencode} de \texttt{forge-init.py} no invoca el adapter (ver
secci\'on~\ref{sec:bug-opencode}).

\subsection{Ausencia de teardown --- resuelto}

\texttt{forge-teardown.py} no exist\'ia en v1. Ahora existe con dry-run por
defecto, eliminaci\'on selectiva de artefactos de forge (no Tier 3 del proyecto),
y 8 tests de cobertura.

\subsection{Ecosistema de profiles insuficiente --- resuelto}

En v1 hab\'ia cuatro profiles. En v2 hay ocho: se agregaron \texttt{fastapi},
\texttt{express}, \texttt{rails}, \texttt{nestjs}. El test param\'etrico en
\texttt{test\_profiles.py} verifica que cada agente de cada profile cumpla el
est\'andar de frontmatter, modelo, secciones y longitud m\'inima.

\subsection{Otras mejoras verificadas}

\begin{itemize}
  \item \textbf{Fases de sprint en CLAUDE.md:} \texttt{\_render\_phases()} lee
    \texttt{sprint.phases} desde \texttt{project.yaml} y genera el listado real.
  \item \textbf{Disclaimer en compliance:} la secci\'on ``Limitaciones'' advierte
    que el agente opera sobre conocimiento de entrenamiento, no texto legal oficial.
  \item \textbf{Compatibilidad Python 3.9:} se usa \texttt{Optional[str]} en lugar
    de \texttt{str | None}.
\end{itemize}

% ----------------------------------------------------------------
\section{Problemas del v1 que persisten en v2}

\subsection{Fix messages del audit no ejecutables}

El audit genera la siguiente cadena como acci\'on correctiva:

\begin{lstlisting}[language=Python, caption={forge-audit.py l\'ineas 223, 242}]
"fix": f"forge-init.py --tool claude-code --force --only={agent['name']}"
\end{lstlisting}

\texttt{forge-init.py} no est\'a en el PATH del sistema. El README muestra el
comando correcto: \texttt{python3 .agentic/scripts/forge-init.py -{}-tool
claude-code}. Un usuario que copie y ejecute la acci\'on correctiva del audit
recibir\'a \texttt{command not found}. Ning\'un test verifica el formato de los
fix messages.

\subsection{Mutaci\'on silenciosa en el hook pre-commit}

El hook sigue haciendo \texttt{git add} sobre \texttt{docs/progress.html} dentro
del proceso de commit, sin que el desarrollador haya inspeccionado ese archivo:

\begin{lstlisting}[language=bash, caption={hooks/pre-commit: mutaci\'on activa}]
if ! git -C "$ROOT" diff --quiet "$HTML"; then
  echo "[forge pre-commit] Actualizando docs/progress.html..."
  git -C "$ROOT" add "$HTML"
fi
\end{lstlisting}

El mensaje informativo agregado en v2 notifica la operaci\'on pero no cambia el
comportamiento. No existe ning\'un test que cubra el hook ni \texttt{token-stats.py}.

\subsection{Similitud de texto como m\'etrica de calidad}

El audit usa \texttt{SequenceMatcher.ratio()} con umbrales fijos
(\texttt{SIMILARITY\_WARN = 0.80}, \texttt{SIMILARITY\_OUTDATED = 0.50}).
Un agente correctamente especializado para Rails puede tener baja similitud con el
core gen\'erico y aparecer como posiblemente desactualizado.

\subsection{Mecanismo de actualizaci\'on destructivo}

\texttt{-{}-force} sobreescribe el agente de destino completamente. No existe
merge, diff interactivo, ni versionado de customizaciones locales.

% ----------------------------------------------------------------
\section{Problemas nuevos detectados en v2}

\subsection{agent-standard.md desincronizado}

La documentaci\'on de referencia de profiles contiene:

\begin{lstlisting}[caption={agent-standard.md: datos incorrectos}]
| `rails`  | `backend-engineer` *(pendiente)* |
| `fastapi` | `backend-engineer` *(pendiente)* |
\end{lstlisting}

Los agentes reales son \texttt{fullstack-engineer} (rails) y
\texttt{api-engineer} (fastapi). Los profiles \texttt{express} y \texttt{nestjs}
no aparecen en la tabla en absoluto.

\subsection{fullstack-engineer sin descripci\'on en forge-init.py}

El profile \texttt{rails} provee \texttt{fullstack-engineer}, pero el diccionario
\texttt{role\_descriptions} en \texttt{forge-init.py} no tiene esa entrada. Un
proyecto que use el profile \texttt{rails} ver\'a ``Agente de implementaci\'on''
(descripci\'on gen\'erica) en su \texttt{AGENTS.md} generado. El adapter de Kiro
s\'i tiene la entrada correcta, creando una inconsistencia entre adapters para el
mismo agente.

\subsection{--tool opencode no invoca el adapter de OpenCode}
\label{sec:bug-opencode}

\texttt{forge-init.py} agrupa \texttt{"claude-code"} y \texttt{"opencode"} bajo el
mismo bloque de c\'odigo:

\begin{lstlisting}[language=Python, caption={forge-init.py l\'ineas 382--389: mismo c\'odigo para ambos tools}]
if tool in ("claude-code", "opencode", "all"):
    label = "Claude Code / OpenCode" if tool == "all" else tool.title()
    init_claude_code(root, forge, config)   # mismo codigo para ambos
    install_claude_commands(root, forge, config)
    init_wiki(root, forge, config)
\end{lstlisting}

Un usuario de OpenCode que ejecute \texttt{forge-init.py -{}-tool opencode}
obtendr\'a una instalaci\'on de Claude Code (\texttt{.claude/agents/}, slash
commands) con exit code 0 y sin mensajes de error. El adapter
\texttt{adapters/opencode/generate-agents-md.py} no se invoca desde
\texttt{forge-init.py} en ning\'un caso.

\subsection{Brechas cr\'iticas en el suite de tests}

El suite no incluye:
\begin{itemize}
  \item Tests para \texttt{token-stats.py} (\'unico script sin cobertura).
  \item Tests para el hook pre-commit.
  \item Tests que validen el formato de los fix messages del audit.
  \item Tests que detecten inconsistencias entre \texttt{agent-standard.md} y el
    c\'odigo.
  \item Tests que ejecuten \texttt{-{}-tool opencode} y verifiquen output
    espec\'ifico de OpenCode.
\end{itemize}

% ----------------------------------------------------------------
\section{Nuevas fortalezas presentes solo en v2}

\subsection{Suite de tests con 112 casos}

La adici\'on m\'as significativa para la madurez del proyecto. Cuatro categor\'ias
complementarias: tests de integraci\'on que ejecutan scripts reales sobre
directorios temporales, tests unitarios con \texttt{importlib} controlando
\texttt{sys.argv}, tests estructurales param\'etricos que recorren todos los
profiles, y tests de regresi\'on con docstrings que documentan los bugs
originales. El \texttt{conftest.py} tiene fixtures con deep merge que permiten
configurar exactamente lo necesario sin boilerplate.

\subsection{Adapters con l\'ogica condicional correcta}

El adapter de Kiro genera \texttt{compliance.md} solo si hay frameworks
configurados en \texttt{project.yaml}. Los tres adapters implementan la misma
l\'ogica de compliance autom\'atico. El adapter de OpenCode lee descriptions desde
el frontmatter de los agentes, respetando la prioridad profiles $>$ core.

\subsection{Test param\'etrico de profiles como garant\'ia de est\'andar}

\texttt{test\_profiles.py} recorre din\'amicamente todos los profiles y verifica
frontmatter completo, modelo correcto, secciones requeridas y longitud m\'inima.
Agregar un profile sin que pase estos tests es imposible sin romper el suite.

\subsection{forge-teardown.py con l\'ogica deliberada}

La distinci\'on entre lo que le pertenece al framework y lo que le pertenece al
proyecto est\'a codificada en la l\'ogica del script. Los agentes Tier 3
sobreviven el teardown.

% ----------------------------------------------------------------
\section{An\'alisis comparativo v1 vs. v2}

\begin{table}[h!]
\centering
\caption{Comparativa de dimensiones clave entre v1 y v2}
\begin{tabular}{@{}lllc@{}}
\toprule
\textbf{Dimensi\'on} & \textbf{v1} & \textbf{v2} & \textbf{Tendencia} \\
\midrule
Profiles implementados & 4 & 8 & $\uparrow$ \\
Tests & 0 & 112 & $\uparrow$ \\
Bug \texttt{install\_agent} & Presente & Resuelto & $\checkmark$ \\
Flag \texttt{-{}-only} & Inexistente & Implementado & $\checkmark$ \\
Adapter OpenCode & Vac\'io & Funcional (no integrado) & $\sim$ \\
Adapter Kiro & Vac\'io & Funcional (integrado) & $\checkmark$ \\
Teardown & Ausente & Presente (parcial) & $\uparrow$ \\
Fix messages ejecutables & No & No & $\times$ \\
Hook pre-commit (mutaci\'on) & Silenciosa & Con notificaci\'on & $\sim$ \\
Disclaimer en compliance & Ausente & Presente & $\checkmark$ \\
\texttt{agent-standard.md} sincronizado & Parcialmente & Menos que v1 & $\downarrow$ \\
\texttt{-{}-tool opencode} invoca adapter & N/A & No & $\times$ \\
\texttt{fullstack-engineer} en role\_descriptions & N/A & Ausente & $\times$ \\
Python 3.9 compatible & No & S\'i & $\checkmark$ \\
\bottomrule
\end{tabular}
\end{table}

% ----------------------------------------------------------------
\section{Conclusiones y recomendaciones}

forge v2 resolvi\'o los problemas m\'as visibles identificados en v1. Este es un
m\'erito real que el an\'alisis debe reconocer. Sin embargo, el patr\'on que
emerge al examinar los problemas nuevos es preocupante: al agregar funcionalidad
(profiles, adapters, scripts), la coherencia interna del framework se degrad\'o.
La documentaci\'on de referencia ahora describe un ecosistema m\'as desactualizado
que en v1. El adapter de OpenCode existe pero est\'a desconectado del flujo
principal. Un agente nuevo (\texttt{fullstack-engineer}) no est\'a registrado en
el script principal.

\subsection{Recomendaciones por perfil de equipo}

\begin{description}
  \item[Adopci\'on inmediata] Equipos de 3--8 personas con Claude Code como
    runtime principal, stack cubierto por los ocho profiles, y disposici\'on a
    complementar la documentaci\'on con lectura del c\'odigo fuente.

  \item[Esperar] Equipos que dependan de la documentaci\'on oficial, usen
    OpenCode como runtime principal, o necesiten que los fix messages del audit
    sean ejecutables sin modificaci\'on.

  \item[Contribuir] Equipos cuyo stack no est\'a en los 8 profiles. El test
    param\'etrico facilita agregar profiles sin overhead adicional.
\end{description}

\subsection{Hoja de ruta recomendada para v3}

\begin{enumerate}
  \item Sincronizar \texttt{agent-standard.md} con el ecosistema real de profiles.
  \item Hacer que \texttt{-{}-tool opencode} invoque el adapter correcto.
  \item Corregir el formato de los fix messages del audit para que sean ejecutables.
  \item Agregar tests que detecten inconsistencias entre capas (documentaci\'on,
    scripts, adapters).
\end{enumerate}

% ----------------------------------------------------------------
\section*{Anexo A: Bugs nuevos encontrados en v2}
\addcontentsline{toc}{section}{Anexo A: Bugs nuevos encontrados en v2}

\subsection*{A.1 Fix messages no ejecutables en el audit}

\textbf{Archivo:} \texttt{scripts/forge-audit.py}, l\'ineas 223 y 242

\begin{lstlisting}[language=Python, caption={Fix message actual --- no ejecutable}]
"fix": f"forge-init.py --tool claude-code --force --only={agent['name']}"
\end{lstlisting}

\texttt{forge-init.py} no est\'a en el PATH del sistema. El comando correcto es
\texttt{python3 .agentic/scripts/forge-init.py -{}-tool claude-code}. Un usuario
que copie y ejecute el fix sugerido recibir\'a \texttt{command not found}.

\textbf{Fix sugerido:}

\begin{lstlisting}[language=Python, caption={Fix message corregido}]
"fix": f"python3 .agentic/scripts/forge-init.py --tool claude-code --force --only={agent['name']}"
\end{lstlisting}

\subsection*{A.2 fullstack-engineer sin entrada en role\_descriptions}

\textbf{Archivo:} \texttt{scripts/forge-init.py}, diccionario
\texttt{role\_descriptions}

El profile \texttt{rails} provee \texttt{fullstack-engineer}, pero ese nombre no
tiene entrada en \texttt{role\_descriptions}. El \texttt{AGENTS.md} generado
mostrar\'a \texttt{"Agente de implementaci\'on"} para ese agente.

\textbf{Fix sugerido:}

\begin{lstlisting}[language=Python, caption={Entrada faltante en role\_descriptions}]
"fullstack-engineer": "Full-stack -- frontend, backend y base de datos en proyectos Rails"
\end{lstlisting}

\subsection*{A.3 --tool opencode no invoca el adapter de OpenCode}

\textbf{Archivo:} \texttt{scripts/forge-init.py}, bloque de selecci\'on de tool
(l\'ineas 382--389)

Ejecutar \texttt{forge-init.py -{}-tool opencode} produce exactamente el mismo
resultado que \texttt{-{}-tool claude-code}. El adapter
\texttt{adapters/opencode/generate-agents-md.py} no se invoca desde
\texttt{forge-init.py} en ning\'un caso. El script termina con exit code 0.

\textbf{Fix sugerido:} agregar un branch espec\'ifico para
\texttt{tool == "opencode"} que invoque el adapter v\'ia subprocess, siguiendo el
mismo patr\'on que \texttt{-{}-tool kiro} usa para invocar
\texttt{adapters/kiro/generate-steering.py}.

\end{document}
